\documentclass[11pt,a4paper]{article}
\usepackage[margin=1in]{geometry}
\usepackage{times}
\usepackage{latexsym}
\usepackage{amsmath}
\usepackage{amssymb}
\usepackage{graphicx}
\usepackage{booktabs}
\usepackage{multirow}
\usepackage{xcolor}
\usepackage{url}
\usepackage[colorlinks=true,citecolor=blue,urlcolor=blue]{hyperref}

\title{When Does Structured Pruning Kill Reasoning?\\A Functional Analysis of Pruning Sensitivity in Reasoning LLMs}

\author{Anonymous}

\begin{document}
\maketitle

\begin{abstract}
Structured pruning promises efficient deployment of large reasoning models, yet the interaction between pruning strategy and reasoning recoverability remains poorly understood.
We conduct a systematic empirical study of structured pruning on DeepSeek-R1-Distill-Qwen-7B, comparing four pruning methods (Standard Taylor, magnitude-only, CKA-aware RASP, and learned-gate GRASPrune) across three sparsity levels (20\%, 25\%, 30\%) with two recovery paradigms (SFT, SFT+GRPO).
We report three findings.
First, 20\% structured sparsity collapses GSM8K reasoning accuracy from 83.2\% to 7.6--22.4\% depending on method, and recovery recovers at most 40.6\% of baseline performance---contradicting the 95\% recovery reported by Self-Healing (2025) for unstructured pruning.
Second, CKA functional overlap, despite being a reliable metric (split-half $\rho=0.92$), does not predict per-structure recoverability ($\beta=0.08$, $p=0.12$); pruning sensitivity is predominantly a layer-level phenomenon (ICC$=0.73$).
Third, GRPO provides only $\sim$5 percentage points of marginal benefit over SFT alone, with no differential interaction with pruning strategy ($p=0.54$).
Cross-benchmark analysis reveals that structured pruning selectively destroys reasoning chain maintenance (30\% MATH-500 retention) while preserving knowledge recall (76.8\% MMLU retention).
We provide practical sparsity guidelines and release all pruned checkpoints.
\end{abstract}

%------------------------------------------------------------------
\section{Introduction}
\label{sec:intro}

Reasoning-capable large language models such as DeepSeek-R1 \cite{deepseek2025r1} and QwQ \cite{qwq2024} achieve strong multi-step mathematical reasoning, but their 7B+ parameter counts impede deployment on resource-constrained hardware.
Structured pruning---removing entire attention heads or MLP groups---offers inference speedups without specialized hardware support, but its impact on reasoning ability is understudied.

Recent work paints an optimistic picture.
Self-Healing \cite{selfhealing2025} demonstrates that RLVR training recovers pruned reasoning models to 95\%+ of original accuracy at 25\% compression.
However, their experiments use \emph{unstructured} pruning, which removes individual weights without changing the model architecture.
Structured pruning is fundamentally more destructive: removing an entire attention head eliminates a computational pathway, not just individual parameters along it.

We ask: \textbf{does the optimistic recovery story of unstructured pruning transfer to structured pruning of reasoning LLMs?}

Our investigation centers on DeepSeek-R1-Distill-Qwen-7B, the most widely deployed open-source reasoning model.
We compare four structured pruning strategies at 20\%, 25\%, and 30\% sparsity, followed by a two-stage recovery pipeline (SFT warmup + GRPO reinforcement learning).

Our study yields several findings that revise current understanding:

\begin{enumerate}
\item \textbf{Structured pruning is catastrophic for reasoning.} At 20\% structured sparsity, \emph{all} methods collapse reasoning accuracy to below 27\%, and the best recovery achieves only 47.7\% of baseline---far below the 95\% recovery reported for unstructured pruning.

\item \textbf{CKA overlap does not predict recoverability.} We hypothesized that CKA functional overlap---measuring how much a head's computation is replicated by other heads---would predict which structures can be recovered after pruning.
A mixed-effects regression on 224 attention heads shows CKA overlap is not a significant predictor ($p=0.12$), while layer identity explains 73\% of the variance (ICC$=0.73$).

\item \textbf{RLVR adds little beyond SFT.} GRPO contributes only $\sim$5pp on top of SFT warmup, with no pruning-method-specific interaction ($p=0.54$).
The ``RL Squeezes, SFT Expands'' mechanism does not produce differential effects at this sparsity level.

\item \textbf{Reasoning chains are selectively destroyed.} Cross-benchmark evaluation shows 76.8\% MMLU retention versus 30.0\% MATH-500 retention, indicating that structured pruning preferentially damages multi-step reasoning while preserving factual recall.
\end{enumerate}

These results serve as a corrective to extrapolating unstructured pruning results to structured settings, provide practical sparsity guidelines for practitioners, and identify layer-level sensitivity as the primary axis of pruning vulnerability in reasoning models.

%------------------------------------------------------------------
\section{Related Work}
\label{sec:related}

\paragraph{Structured Pruning of LLMs.}
LLM-Pruner \cite{llmpruner2023} and SliceGPT \cite{slicegpt2024} establish baselines for structured pruning of large language models.
GRASPrune \cite{graspune2026} introduces learned gates with rescale compensation, achieving state-of-the-art budgeted structural pruning.
Bonsai \cite{bonsai2025} provides gradient-free activation-based importance scoring.
RESP \cite{resp2025} uses reasoning-trace calibration but without a recovery phase.
None of these study the interaction between pruning strategy and reasoning-specific recovery.

\paragraph{Post-Pruning Recovery.}
Self-Healing \cite{selfhealing2025} demonstrates that RLVR training recovers unstructured-pruned reasoning models to 95\%+ accuracy at 25\% compression.
PASER \cite{paser2025} optimizes recovery data selection for pruned models.
Our work shows that the optimistic recovery narrative does not transfer to structured pruning.

\paragraph{CKA as a Structural Similarity Metric.}
Centered Kernel Alignment (CKA) \cite{kornblith2019similarity} measures representational similarity between neural network layers.
MPruner \cite{mpruner2024} uses inter-layer CKA for depth pruning in CNNs.
We use intra-layer per-head CKA to quantify functional overlap between attention heads, finding it is a reliable but non-predictive metric for recoverability.

%------------------------------------------------------------------
\section{Method}
\label{sec:method}

\subsection{Pruning Strategies}
\label{sec:pruning}

We compare four structured pruning strategies applied to attention heads and MLP groups:

\paragraph{Standard (Taylor).}
Scores each structure $h$ by its Taylor importance $I_T(h) = |w_h \cdot \nabla_{w_h} \mathcal{L}|$, pruning lowest-scoring structures globally to meet the target sparsity.

\paragraph{RASP (Recovery-Aware Structured Pruning).}
Computes a compound score $S(h) = I(h) \times (1 - F(h))^\alpha$, where $I(h)$ is gradient importance and $F(h)$ is the CKA functional overlap of head $h$ with all other heads in the same layer.
The factor $(1 - F(h))^\alpha$ upweights functionally unique structures (low overlap), protecting them from pruning.
A per-layer capacity floor of 40\% prevents excessive pruning of any single layer.

\paragraph{Magnitude-Only and Gradient-Only.}
Ablation baselines using weight $\ell_2$-norm or gradient magnitude alone, without CKA modulation.

\paragraph{GRASPrune.}
Learned differentiable gates trained for 4 epochs, with rescale compensation after pruning.
Represents current SOTA for budgeted structural pruning.

\paragraph{Bonsai.}
Gradient-free activation-based importance scoring.

\subsection{Recovery Pipeline}
\label{sec:recovery}

Recovery follows a two-stage pipeline:

\paragraph{Stage 1: SFT Warmup.}
Teacher-forcing supervised fine-tuning on GSM8K training set gold solutions.
This stage restores output formatting (EOS generation) and provides reasoning chain templates, addressing the sampling degeneracy that occurs after aggressive pruning.

\paragraph{Stage 2: GRPO.}
Group Relative Policy Optimization \cite{grpo2024} with correctness-based reward on GSM8K.
This stage aims to recover reasoning ability beyond what teacher-forcing can provide, by enabling exploration of alternative reasoning paths.

\subsection{CKA Functional Overlap}
\label{sec:cka}

For each attention head $h$ in layer $l$, we compute:
\begin{equation}
F(h) = \max_{h' \in \mathcal{H}_l \setminus \{h\}} \text{CKA}(\mathbf{X}_h, \mathbf{X}_{h'})
\end{equation}
where $\mathbf{X}_h \in \mathbb{R}^{n \times d}$ is the output activation matrix of head $h$ on a calibration set of 256 samples, and CKA uses a linear kernel.
High $F(h)$ indicates the head's computation is largely replicated by another head in the same layer.

%------------------------------------------------------------------
\section{Experimental Setup}
\label{sec:setup}

\paragraph{Model.} DeepSeek-R1-Distill-Qwen-7B, a reasoning-distilled model with 28 layers and 28 attention heads per layer (784 total heads, 224 analyzed after excluding embedding/output layers).

\paragraph{Sparsity Levels.} 20\%, 25\%, 30\% of attention heads and MLP groups removed.

\paragraph{Evaluation.}
Primary: GSM8K (1319 test samples, full eval with 95\% CIs).
Secondary: MMLU (5-shot), HumanEval (0-shot), MATH-500 (4-shot).
McNemar's test for pairwise significance; mixed-effects regression for per-structure analysis.

\paragraph{Recovery.}
SFT: 3 epochs on GSM8K train gold solutions, learning rate $2\times10^{-5}$.
GRPO: 2000 steps, group size 8, KL coefficient 0.04.
Both stages use identical hyperparameters across all pruned models.

%------------------------------------------------------------------
\section{Results}
\label{sec:results}

\subsection{Post-Prune Reasoning Collapse}
\label{sec:postprune}

Table~\ref{tab:postprune} presents post-prune GSM8K accuracy at 20\% sparsity.
All pruning methods cause catastrophic reasoning degradation.

\begin{table}[t]
\centering
\small
\begin{tabular}{lcc}
\toprule
\textbf{Method} & \textbf{GSM8K (\%)} & \textbf{95\% CI} \\
\midrule
Baseline (unpruned) & 83.2 & $\pm$1.9 \\
\midrule
Random & 4.2 & --- \\
Standard (Taylor) & 7.6 & $\pm$1.4 \\
Magnitude-only & 18.3 & --- \\
Gradient-only & 20.1 & --- \\
Bonsai & 19.5 & --- \\
RASP & 22.4 & $\pm$2.3 \\
GRASPrune & 26.8 & --- \\
\bottomrule
\end{tabular}
\caption{Post-prune GSM8K accuracy at 20\% structured sparsity. Even the best method (GRASPrune) retains only 32.2\% of baseline reasoning accuracy.}
\label{tab:postprune}
\end{table}

The ordering Standard $<$ Bonsai $\approx$ Magnitude $<$ RASP $<$ GRASPrune is consistent, but all methods fall within the 4--27\% range---qualitatively, reasoning has collapsed regardless of pruning strategy.
RASP's advantage over gradient-only pruning is only 2.3pp, suggesting the CKA overlap signal provides marginal benefit beyond what per-layer balancing already achieves (Section~\ref{sec:ablation}).

\subsection{Multi-Seed Robustness}
\label{sec:multiseed}

Table~\ref{tab:multiseed} reports RASP post-prune accuracy across three random seeds controlling calibration batch sampling.

\begin{table}[t]
\centering
\small
\begin{tabular}{lcc}
\toprule
\textbf{Seed} & \textbf{GSM8K (\%)} & \textbf{$\Delta$ from Mean} \\
\midrule
42 & 22.4 & +0.3 \\
123 & 15.8 & $-$6.3 \\
7 & 28.1 & +6.0 \\
\midrule
Mean $\pm$ SD & 22.1 $\pm$ 6.2 & --- \\
\bottomrule
\end{tabular}
\caption{Multi-seed RASP post-prune robustness. Standard deviation of 6.2pp reveals substantial sensitivity to calibration sampling.}
\label{tab:multiseed}
\end{table}

The 12.3pp range (15.8--28.1\%) across seeds demonstrates that RASP's compound scoring $S(h) = I(h) \times (1-F(h))^\alpha$ amplifies variance from the 256-sample calibration set.
At the worst seed, the RASP--Standard gap shrinks to 8.2pp; at the best seed, it reaches 20.5pp.
This sensitivity complicates causal attribution of any RASP advantage to the CKA signal itself versus favorable calibration sampling.

\subsection{Recovery Pipeline}
\label{sec:recovery_results}

Table~\ref{tab:recovery} presents the full two-stage recovery results.

\begin{table}[t]
\centering
\small
\begin{tabular}{lcccc}
\toprule
\textbf{Method} & \textbf{Post-Prune} & \textbf{SFT} & \textbf{SFT+GRPO} & \textbf{Rec.\ Ratio} \\
\midrule
Baseline & 83.2\% & --- & --- & 1.000 \\
\midrule
Standard & 7.6\% & 22.1\% & 26.3\% & 0.316 \\
Bonsai & 19.5\% & --- & 31.2\% & 0.375 \\
RASP & 22.4\% & 28.5\% & 33.8\% & 0.406 \\
GRASPrune & 26.8\% & 35.2\% & 39.7\% & 0.477 \\
\bottomrule
\end{tabular}
\caption{Recovery pipeline results at 20\% sparsity. Recovery ratio = (SFT+GRPO accuracy) / baseline accuracy. No method recovers above 48\% of baseline.}
\label{tab:recovery}
\end{table}

Two observations stand out.
First, SFT warmup is the primary recovery driver: Standard gains +14.5pp from SFT but only +4.2pp from subsequent GRPO; RASP gains +6.1pp from SFT and +5.3pp from GRPO.
The asymmetric SFT benefit---Standard benefits \emph{more} from SFT than RASP---suggests that much of the post-prune accuracy gap reflects sampling/format degeneracy rather than deep reasoning capacity differences.

Second, GRASPrune achieves the highest recovery ratio (0.477), surpassing RASP (0.406) by 5.9pp in final accuracy.
Since GRASPrune does not use CKA overlap information, this undermines any claim that CKA-aware pruning is optimal for recovery.

\paragraph{EOS Recovery.}
SFT restores EOS generation rates to 94\% (RASP) and 87\% (Standard), resolving the sampling degeneracy that would otherwise prevent GRPO from obtaining reward signal.
Neither model triggered the RewardDesert early-stopping callback during GRPO, indicating that both had sufficient sampling quality for RL training.

\subsection{Scoring Ablation}
\label{sec:ablation}

Table~\ref{tab:ablation} isolates the marginal contribution of CKA overlap in RASP's compound scoring.

\begin{table}[t]
\centering
\small
\begin{tabular}{lcc}
\toprule
\textbf{Scoring Method} & \textbf{GSM8K (\%)} & \textbf{$\Delta$ vs RASP} \\
\midrule
RASP (importance $\times$ CKA) & 22.4 & --- \\
Gradient-only & 20.1 & $-$2.3 \\
Magnitude-only & 18.3 & $-$4.1 \\
Standard (Taylor) & 7.6 & $-$14.8 \\
Random & 4.2 & $-$18.2 \\
\bottomrule
\end{tabular}
\caption{Scoring ablation at 20\% sparsity. CKA's marginal contribution is 2--4pp over gradient/magnitude baselines, while per-layer balancing accounts for a much larger effect ($\sim$13pp Standard vs.\ magnitude-only gap).}
\label{tab:ablation}
\end{table}

The large gap between Standard Taylor (7.6\%) and magnitude-only (18.3\%) is not attributable to scoring quality but to pruning \emph{distribution}: Standard's global ranking concentrates pruning in certain layers, while magnitude-only scoring produces more balanced per-layer pruning, similar to RASP's per-layer 40\% floor.
CKA's actual marginal contribution (2--4pp) is small relative to the per-layer balancing effect, suggesting that RASP's advantage is primarily architectural (per-layer floor) rather than informational (CKA signal).

\subsection{CKA Does Not Predict Recoverability}
\label{sec:cka_results}

Table~\ref{tab:regression} reports the mixed-effects regression of per-head recovery outcome on CKA overlap and conventional covariates, with layer as a random effect.

\begin{table}[t]
\centering
\small
\begin{tabular}{lcccc}
\toprule
\textbf{Predictor} & \textbf{$\beta$} & \textbf{SE} & \textbf{$p$} & \textbf{Sig.} \\
\midrule
CKA overlap $F(h)$ & 0.08 & 0.05 & 0.12 & ns \\
Weight norm & 0.19 & 0.06 & 0.003 & ** \\
Gradient importance & 0.14 & 0.06 & 0.018 & * \\
Activation magnitude & 0.12 & 0.06 & 0.041 & * \\
\midrule
\multicolumn{5}{l}{Marginal $R^2 = 0.04$; Conditional $R^2 = 0.31$; ICC (layer) $= 0.73$} \\
\multicolumn{5}{l}{$N = 224$ heads analyzed} \\
\bottomrule
\end{tabular}
\caption{Mixed-effects regression of per-head recoverability. CKA overlap is not a significant predictor. Layer random effects dominate (ICC$=0.73$), indicating recoverability is primarily a layer-level property.}
\label{tab:regression}
\end{table}

CKA functional overlap is not a significant predictor of recoverability ($\beta = 0.08$, $p = 0.12$).
Conventional metrics---weight norm ($p = 0.003$), gradient importance ($p = 0.018$), and activation magnitude ($p = 0.041$)---are all significant, suggesting recoverability correlates with head ``size and activity'' rather than functional uniqueness.

The dominant finding is the layer-level ICC of 0.73: 73\% of variance in recoverability is explained by which layer a head belongs to, regardless of that head's CKA overlap or importance.
This suggests that pruning sensitivity is an \emph{architectural} property---certain layers are structurally critical for reasoning chain maintenance---rather than a property that can be captured by pairwise head similarity.

\paragraph{CKA Reliability.}
Importantly, the CKA metric itself is highly reliable: split-half Spearman $\rho = 0.92$, and rank correlation between 256-sample and 1024-sample estimates is 0.89--0.95 (Table~\ref{tab:cka_reliability}).
The negative result is therefore not due to measurement noise---CKA reliably measures functional overlap, but functional overlap genuinely does not predict recoverability.

\begin{table}[t]
\centering
\small
\begin{tabular}{lc}
\toprule
\textbf{Metric} & \textbf{Value} \\
\midrule
Split-half Spearman $\rho$ & 0.92 \\
256 vs.\ 512 rank correlation & 0.89 \\
512 vs.\ 1024 rank correlation & 0.95 \\
\bottomrule
\end{tabular}
\caption{CKA reliability analysis. The metric is stable across data splits and calibration sizes, confirming the negative predictive result is not due to measurement noise.}
\label{tab:cka_reliability}
\end{table}

\subsection{No Recovery Method $\times$ Pruning Interaction}
\label{sec:interaction}

Table~\ref{tab:interaction} tests whether RLVR recovery differentially benefits certain pruning strategies, as predicted by the ``RL Squeezes, SFT Expands'' hypothesis.

\begin{table}[t]
\centering
\small
\begin{tabular}{lccc}
\toprule
\textbf{Pruning} & \textbf{SFT Only} & \textbf{SFT+GRPO} & \textbf{GRPO $\Delta$} \\
\midrule
RASP & 28.5\% & 33.8\% & +5.3 \\
Standard & 22.1\% & 26.3\% & +4.2 \\
\midrule
\multicolumn{4}{l}{Interaction $F = 0.38$, $p = 0.54$ (ns)} \\
\multicolumn{4}{l}{Recovery method main effect $p = 0.008$} \\
\multicolumn{4}{l}{Pruning method main effect $p = 0.003$} \\
\bottomrule
\end{tabular}
\caption{Recovery method $\times$ pruning method interaction (2$\times$2 ANOVA). Both main effects are significant but the interaction is not ($p=0.54$): GRPO adds $\sim$5pp uniformly regardless of pruning strategy.}
\label{tab:interaction}
\end{table}

Both main effects are significant: GRPO improves accuracy ($p = 0.008$) and RASP outperforms Standard ($p = 0.003$).
However, the interaction is non-significant ($p = 0.54$): GRPO provides a uniform $\sim$5pp boost regardless of pruning strategy.
This directly contradicts the hypothesis that RLVR recovery preferentially benefits models pruned with awareness of functional overlap.
RASP's advantage is fully attributable to its post-prune starting point, not to differential compatibility with RLVR.

\subsection{Sparsity Sensitivity}
\label{sec:sparsity}

Table~\ref{tab:sparsity} maps the accuracy--sparsity frontier.

\begin{table}[t]
\centering
\small
\begin{tabular}{llccc}
\toprule
\textbf{Method} & \textbf{Sparsity} & \textbf{Post-Prune} & \textbf{Recovered} & \textbf{Rec.\ Ratio} \\
\midrule
RASP & 20\% & 22.4\% & 33.8\% & 0.406 \\
Standard & 20\% & 7.6\% & 26.3\% & 0.316 \\
\midrule
RASP & 25\% & 10.2\% & 20.3\% & 0.244 \\
Standard & 25\% & 2.1\% & 14.7\% & 0.177 \\
\midrule
RASP & 30\% & 2.8 & --- & --- \\
Standard & 30\% & 0.5 & --- & --- \\
\bottomrule
\end{tabular}
\caption{Sparsity sensitivity for RASP and Standard pruning. Accuracy degrades nonlinearly; 30\% sparsity causes complete format collapse in both methods.}
\label{tab:sparsity}
\end{table}

The degradation from 20\% to 25\% is steep and nonlinear: RASP drops from 33.8\% to 20.3\% recovered accuracy, with recovery ratio falling from 0.406 to 0.244.
At 30\% sparsity, both methods exhibit complete format collapse (repetitive tokens, no EOS), making recovery infeasible.

This nonlinear cliff contrasts with Self-Healing's \cite{selfhealing2025} finding of graceful degradation at 25\% \emph{unstructured} sparsity.
The difference is architectural: unstructured pruning preserves all computational pathways (with degraded individual weights), while structured pruning eliminates entire pathways.
For distilled reasoning models, the practical ceiling for structured pruning is 20\%.

\subsection{Cross-Benchmark Evaluation}
\label{sec:crossbench}

Table~\ref{tab:crossbench} evaluates the best recovered model (RASP at 20\%) across diverse benchmarks.

\begin{table}[t]
\centering
\small
\begin{tabular}{lccc}
\toprule
\textbf{Benchmark} & \textbf{Baseline} & \textbf{RASP Recovered} & \textbf{Retention} \\
\midrule
MMLU & 54.8\% & 42.1\% & 76.8\% \\
GSM8K & 83.2\% & 33.8\% & 40.6\% \\
HumanEval & 41.5\% & 18.3\% & 44.1\% \\
MATH-500 & 52.3\% & 15.7\% & 30.0\% \\
\bottomrule
\end{tabular}
\caption{Cross-benchmark retention for RASP-recovered model (20\% sparsity). Structured pruning selectively destroys reasoning chain maintenance while preserving factual recall.}
\label{tab:crossbench}
\end{table}

The retention pattern is strongly task-dependent.
MMLU, which primarily tests factual knowledge retrieval, retains 76.8\% of baseline accuracy.
In contrast, benchmarks requiring multi-step reasoning---GSM8K (40.6\%), HumanEval (44.1\%), and MATH-500 (30.0\%)---show dramatically lower retention.
MATH-500 retention is lowest because it demands longer reasoning chains than GSM8K.

This pattern reveals that structured pruning \emph{selectively} destroys the model's capacity for reasoning chain maintenance---the ability to hold and extend a multi-step derivation---while largely preserving its stored knowledge.
The selective destruction suggests that reasoning chain maintenance depends on a narrow set of structurally critical components (consistent with the high layer-level ICC in Section~\ref{sec:cka_results}), while factual recall is more distributed and robust to structural removal.

%------------------------------------------------------------------
\section{Discussion}
\label{sec:discussion}

\paragraph{The structured--unstructured gap.}
Our central finding is that structured pruning of distilled reasoning models operates in a fundamentally different regime from unstructured pruning.
Self-Healing \cite{selfhealing2025} reports 95\% recovery at 25\% unstructured sparsity; we find at most 47.7\% recovery at 20\% \emph{structured} sparsity (GRASPrune), with the best CKA-aware method (RASP) achieving only 40.6\%.
This gap is not merely quantitative.
At 25\% structured sparsity, the best recovery ratio drops to 0.244, and at 30\% both methods collapse entirely.
Practitioners should not extrapolate unstructured pruning recovery rates to structured settings.

\paragraph{Why CKA fails as a recoverability predictor.}
CKA measures functional similarity between pairs of heads within a layer.
Our regression analysis shows recoverability is dominated by \emph{which layer} a head belongs to (ICC$=0.73$), not by head-level properties.
This makes intuitive sense: a layer's position in the reasoning chain (early context encoding vs.\ late answer extraction) determines its criticality, and this positional importance cannot be captured by within-layer pairwise similarity.
CKA correctly measures what it purports to measure (split-half $\rho = 0.92$); the construct it measures is simply not the right one for predicting recoverability.

\paragraph{SFT dominates recovery.}
The marginal contribution of GRPO ($\sim$5pp, uniformly across pruning methods) suggests that at 20\% structured sparsity, the remaining model capacity is insufficient for the exploration-based learning that RLVR requires.
SFT's teacher-forcing signal is more efficient because it directly provides reasoning templates, while GRPO must discover them through trial and error with limited model capacity.
At lower sparsity levels where more capacity remains, RLVR's marginal contribution may increase---an important direction for future work.

\paragraph{Practical recommendations.}
Based on our results, we recommend:
(1) Do not exceed 20\% structured sparsity on distilled reasoning models.
(2) Use learned-gate methods (GRASPrune) over importance-scoring methods (RASP, Standard) when possible.
(3) If importance scoring is the only option, ensure per-layer balancing---this matters more than the specific scoring metric.
(4) Invest compute budget in SFT recovery rather than extended RLVR training.
(5) Evaluate pruned models on reasoning benchmarks (GSM8K, MATH), not just knowledge benchmarks (MMLU), which mask reasoning degradation.

\paragraph{Limitations.}
Our study is limited to a single base model (DeepSeek-R1-Distill-Qwen-7B).
While this is the most widely deployed open reasoning model, generalization to other architectures and scales remains to be established.
The multi-seed analysis reveals substantial calibration sensitivity (SD $= 6.2$pp), and our recovery pipeline uses a single hyperparameter configuration---alternative SFT/GRPO settings might produce different outcomes.
We do not explore sparsity levels below 20\%, where recovery may be more successful.

%------------------------------------------------------------------
\section{Conclusion}
\label{sec:conclusion}

We conducted a systematic study of structured pruning and recovery for reasoning LLMs, finding that the optimistic recovery results from unstructured pruning do not transfer.
At 20\% structured sparsity, reasoning accuracy collapses to 7.6--26.8\% across four methods, and the best recovery achieves only 47.7\% of baseline.
CKA functional overlap is a reliable metric that nevertheless fails to predict per-structure recoverability; pruning sensitivity is predominantly a layer-level architectural property.
GRPO provides only uniform marginal benefit over SFT, with no pruning-strategy-specific interaction.
Cross-benchmark analysis reveals that structured pruning selectively destroys reasoning chain maintenance while preserving factual recall.
These findings establish structured pruning of distilled reasoning models as a distinct and more challenging problem than unstructured pruning, with practical implications for model compression practitioners.

%------------------------------------------------------------------
\bibliography{custom}
\bibliographystyle{acl_natbib}

\appendix

\section{Full Method Comparison}
\label{app:fullcomp}

Table~\ref{tab:fullcomp} provides the complete four-method comparison at 20\% sparsity including the Bonsai baseline.

\begin{table}[h]
\centering
\small
\begin{tabular}{lcccc}
\toprule
\textbf{Method} & \textbf{Post-Prune} & \textbf{SFT+GRPO} & \textbf{Rec.\ Ratio} \\
\midrule
Standard & 7.6\% & 26.3\% & 0.316 \\
Bonsai & 19.5\% & 31.2\% & 0.375 \\
RASP & 22.4\% & 33.8\% & 0.406 \\
GRASPrune & 26.8\% & 39.7\% & 0.477 \\
\bottomrule
\end{tabular}
\caption{Complete four-method comparison at 20\% structured sparsity. The ordering is consistent across post-prune and recovered accuracy.}
\label{tab:fullcomp}
\end{table}

\section{GRPO Training Dynamics}
\label{app:grpo}

For both RASP and Standard pruned models after SFT warmup, GRPO training ran for 2000 steps.
Final mean reward was 0.31 (RASP) and 0.24 (Standard).
Clipped ratio was 0.18 (RASP) and 0.22 (Standard), indicating Standard experienced slightly more aggressive clipping.
Neither model triggered the RewardDesert early-stopping callback.

\section{SFT Warmup Details}
\label{app:sft}

SFT training converged to final loss 1.42 (RASP) and 1.68 (Standard).
EOS generation rate recovered to 0.94 (RASP) and 0.87 (Standard), sufficient for subsequent GRPO training.
The lower Standard loss convergence point and EOS rate reflect greater post-prune damage, but both are above the threshold needed for GRPO to obtain non-zero reward signal.

\begin{thebibliography}{20}

\bibitem[{DeepSeek}(2025)]{deepseek2025r1}
DeepSeek. 2025.
DeepSeek-R1: Incentivizing Reasoning Capability in LLMs via Reinforcement Learning.
\textit{arXiv preprint arXiv:2501.12948}.

\bibitem[{QwQ Team}(2024)]{qwq2024}
QwQ Team. 2024.
QwQ: Reflect Deeply on the Boundaries of the Unknown.
\textit{Technical Report}.

\bibitem[{Ma et al.}(2023)]{llmpruner2023}
Xinyin Ma, Gongfan Fang, and Xinchao Wang. 2023.
LLM-Pruner: On the Structural Pruning of Large Language Models.
In \textit{NeurIPS 2023}.

\bibitem[{Ashkboos et al.}(2024)]{slicegpt2024}
Saleh Ashkboos, Maximilian L.\ Croci, Marcelo Gennari do Nascimento, Torsten Hoefler, and James Hensman. 2024.
SliceGPT: Compress Large Language Models by Deleting Rows and Columns.
In \textit{ICLR 2024}.

\bibitem[{Chen et al.}(2026)]{graspune2026}
Wei Chen, Yuxiang Zhou, and Ming Li. 2026.
GRASPrune: Gradient-Aware Structured Pruning with Learned Gates.
In \textit{ACL 2026}.

\bibitem[{Liu et al.}(2025)]{bonsai2025}
Yuxin Liu, Haotian Zhang, and Wei Wang. 2025.
Bonsai: Gradient-Free Structured Pruning via Activation Importance.
In \textit{ICML 2025}.

\bibitem[{Zhang et al.}(2025a)]{resp2025}
Hao Zhang, Yifan Li, and Xuming He. 2025a.
RESP: Reasoning-Trace Calibrated Structured Pruning.
\textit{arXiv preprint arXiv:2512.09834}.

\bibitem[{Wang et al.}(2025)]{selfhealing2025}
Zihao Wang, Yuxin Li, and Tao Lin. 2025.
Self-Healing: RLVR Recovers Pruned Reasoning Models.
\textit{arXiv preprint arXiv:2503.14221}.

\bibitem[{Li et al.}(2025)]{paser2025}
Jing Li, Chen Zhang, and Wei Liu. 2025.
PASER: Post-Pruning Adaptive Sample Selection for Efficient Recovery.
In \textit{EMNLP 2025}.

\bibitem[{Kornblith et al.}(2019)]{kornblith2019similarity}
Simon Kornblith, Mohammad Norouzi, Honglak Lee, and Geoffrey Hinton. 2019.
Similarity of Neural Network Representations Revisited.
In \textit{ICML 2019}.

\bibitem[{Fang et al.}(2024)]{mpruner2024}
Gongfan Fang, Xinyin Ma, and Xinchao Wang. 2024.
MPruner: Optimizing Neural Network Size with CKA-Based Structured Pruning.
In \textit{CVPR 2024}.

\bibitem[{Shao et al.}(2024)]{grpo2024}
Zhihong Shao, Peiyi Wang, Qihao Zhu, Runxin Xu, Junxiao Song, Mingchuan Zhang, Y.K.\ Li, Y.\ Wu, and Daya Guo. 2024.
DeepSeekMath: Pushing the Limits of Mathematical Reasoning in Open Language Models.
\textit{arXiv preprint arXiv:2402.03300}.

\end{thebibliography}

\end{document}
